\documentclass[11pt]{article}
\usepackage{amsmath,amssymb}
\usepackage{geometry}
\usepackage{setspace}
\usepackage{enumitem}
\geometry{margin=1in}
\setstretch{1.15}

\begin{document}

\begin{center}
{\Large \textbf{CSE 400 — Lecture 15 Scribe}}\\
\vspace{0.2em}
Fundamentals of Probability in Computing\\
Joint Random Variables and Joint Distributions
\end{center}

\section{Motivation: Why Study Joint Random Variables}

In many real-world systems, multiple random variables must be analyzed simultaneously.  
The behavior of one variable often depends on another.

\subsection*{Example 1: Smart Transportation}

Random variables:
\begin{itemize}[noitemsep]
\item $X$ = Vehicle Speed
\item $Y$ = Traffic Density
\end{itemize}

Question:

Can speed be high when traffic density is also high?

This indicates that the two variables may be correlated random variables.

Applications:
\begin{itemize}[noitemsep]
\item Adaptive traffic signals
\item Congestion prediction
\item Autonomous vehicle routing
\end{itemize}

\subsection*{Example 2: Wireless Network Performance}

Random variables:
\begin{itemize}[noitemsep]
\item $X$ = Signal-to-Noise Ratio (SNR)
\item $Y$ = Data Throughput
\end{itemize}

Question:

If SNR is high, will throughput always be high?

Throughput depends on SNR, but may also depend on other conditions such as network congestion.

Applications:
\begin{itemize}[noitemsep]
\item 5G scheduling
\item Adaptive modulation
\item Network planning
\end{itemize}

\subsection*{Additional Examples}

\begin{itemize}[noitemsep]
\item Weather prediction: $X$ = Rainfall, $Y$ = Temperature
\item Drone delivery: $X$ = Wind Speed, $Y$ = Delivery Time
\item Rolling two dice
\end{itemize}

These situations require joint probability calculations.

\section{Discrete Case — Joint PMF}

\subsection{Setup}

Suppose $X$ and $Y$ are discrete random variables.

\[
X \in \{x_1, x_2, \ldots, x_n\}
\]

\[
Y \in \{y_1, y_2, \ldots, y_m\}
\]

The ordered pair $(X,Y)$ takes values in the product set

\[
\{(x_1,y_1),(x_1,y_2),\ldots,(x_n,y_m)\}
\]

\subsection{Definition: Joint PMF}

The joint probability mass function is defined as

\[
p(x_i,y_j) = P(X=x_i, Y=y_j)
\]

This represents the probability that the two outcomes occur simultaneously.

\section{Joint PMF Table Representation}

The joint PMF can be organized in a table.

\[
\begin{array}{c|cccc}
X\backslash Y & y_1 & y_2 & \cdots & y_m \\ \hline
x_1 & p(x_1,y_1) & p(x_1,y_2) & \cdots & p(x_1,y_m) \\
x_2 & p(x_2,y_1) & p(x_2,y_2) & \cdots & p(x_2,y_m) \\
\vdots & \vdots & \vdots & & \vdots \\
x_n & p(x_n,y_1) & p(x_n,y_2) & \cdots & p(x_n,y_m)
\end{array}
\]

Key properties:

\[
0 \le p(x_i,y_j) \le 1
\]

\[
\sum_{i=1}^{n}\sum_{j=1}^{m} p(x_i,y_j) = 1
\]

\section{Example: Rolling Two Dice}

\subsection*{Experiment}

Roll two fair dice.

Let
\begin{itemize}[noitemsep]
\item $X$ = value on the first die
\item $T$ = total (sum) of both dice
\end{itemize}

Each outcome $(i,j)$ has probability

\[
\frac{1}{36}
\]

Construct the joint probability table of $(X,T)$.

\subsection*{Joint Probability Table}

\[
\begin{array}{c|cccccc}
X\backslash T & 2 & 3 & 4 & 5 & 6 & 7 \\ \hline
1 & \frac{1}{36} & \frac{1}{36} & \frac{1}{36} & \frac{1}{36} & \frac{1}{36} & \frac{1}{36} \\
2 & 0 & \frac{1}{36} & \frac{1}{36} & \frac{1}{36} & \frac{1}{36} & \frac{1}{36} \\
3 & 0 & 0 & \frac{1}{36} & \frac{1}{36} & \frac{1}{36} & \frac{1}{36} \\
4 & 0 & 0 & 0 & \frac{1}{36} & \frac{1}{36} & \frac{1}{36} \\
5 & 0 & 0 & 0 & 0 & \frac{1}{36} & \frac{1}{36} \\
6 & 0 & 0 & 0 & 0 & 0 & \frac{1}{36}
\end{array}
\]

Observations:
\begin{itemize}[noitemsep]
\item Entries are either $0$ or $\frac{1}{36}$.
\item $X$ and $T$ are not independent.
\end{itemize}

\section{Continuous Case — Joint PDF}

The continuous case is analogous to the discrete case.

Correspondence:

\[
\begin{array}{c|c}
\text{Discrete} & \text{Continuous} \\
\hline
\text{Values} & \text{Intervals} \\
\text{Joint PMF} & \text{Joint PDF} \\
\text{Sums} & \text{Double Integrals}
\end{array}
\]

If

\[
X \in [a,b], \qquad Y \in [c,d]
\]

then the pair $(X,Y)$ lies in the product region

\[
[a,b] \times [c,d]
\]

\subsection{Definition: Joint PDF}

A function $f(x,y)$ is called the joint probability density function of $X$ and $Y$ if

\[
P((X,Y)\ \text{lies in a small rectangle of area}\ dx\,dy)
= f(x,y)\,dx\,dy
\]

Key properties:

\[
f(x,y) \ge 0
\]

\[
\int_c^d \int_a^b f(x,y)\,dx\,dy = 1
\]

Note:

The joint PDF $f(x,y)$ may take values greater than $1$, because it is a density, not a probability.

\section{Worked Example}

Given

\[
f(x,y) = 4xy
\]

for

\[
0 < x < 1, \qquad 0 < y < 1
\]

Thus $(X,Y)$ lies in the unit square.

Tasks:
\begin{enumerate}[label=\arabic*.]
\item Show that $f(x,y)$ is a valid joint PDF.
\item Visualize the event.
\item Compute $P(A)$.
\end{enumerate}

\subsection*{Example Solution}

\[
A = \{X < 0.5,\ Y > 0.5\}
\]

\subsubsection*{Step 1: Check Validity (Normalization)}

Compute

\[
\int_0^1 \int_0^1 4xy\,dx\,dy
\]

First integrate with respect to $x$:

\[
\int_0^1 \left[\int_0^1 4xy\,dx\right] dy
\]

\[
= \int_0^1 2y\,dy
\]

Thus

\[
f(x,y) \ge 0
\]

and total probability equals $1$.

Therefore $f(x,y)$ is a valid joint PDF.

\subsubsection*{Step 2: Compute $P(A)$}

\[
A = \{X < 0.5,\ Y > 0.5\}
\]

\[
P(A) = \int_0^{0.5} \int_{0.5}^{1} 4xy\,dy\,dx
\]

First integrate with respect to $y$:

\[
\int_0^{0.5} \left[2xy^2\right]_{0.5}^{1} dx
\]

\[
= \int_0^{0.5} 2x(1 - 0.25) dx
\]

\[
= \frac{3}{16}
\]

\section{Joint Cumulative Distribution Function (Joint CDF)}

Suppose $X$ and $Y$ are jointly distributed random variables.

Notation:

\[
\{X \le x,\ Y \le y\}
\]

\subsection*{Definition}

The joint cumulative distribution function is

\[
F(x,y) = P(X \le x,\ Y \le y)
\]

\section{Joint CDF — Continuous Case}

If $X$ and $Y$ have joint density $f(x,y)$ over

\[
[a,b] \times [c,d]
\]

then

\[
F(x,y) = \int_{-\infty}^{x} \int_{-\infty}^{y} f(u,v)\,du\,dv
\]

\subsection*{Recovering the Joint PDF}

Since two variables are involved, use partial derivatives:

\[
f(x,y) = \frac{\partial^2}{\partial x \partial y} F(x,y)
\]

\section{Joint CDF — Discrete Case}

If $X$ and $Y$ are discrete with joint PMF $p(x_i,y_j)$,

\[
F(x,y) = \sum_{x_i \le x} \sum_{y_j \le y} p(x_i,y_j)
\]

Interpretation:

To compute $F(x,y)$:
\begin{itemize}[noitemsep]
\item Add probabilities in the table
\item Include rows where $x_i \le x$
\item Include columns where $y_j \le y$
\end{itemize}

\section{Properties of the Joint CDF}

\subsection*{1. Non-decreasing}

$F(x,y)$ is non-decreasing.

If $x$ or $y$ increases, $F(x,y)$ must stay the same or increase.

\subsection*{2. Lower-Left Boundary}

\[
F(x,y) = 0
\]

at the lower-left of the joint range.

\subsection*{3. Upper-Right Boundary}

\[
F(x,y) = 1
\]

at the upper-right of the joint range.

\end{document}